\manualchapter{Patches and troubleshooting}{sec:operation}
\subsection{Clocked percussion and coupled voices}
Patch a short envelope to Tone 1 Level and sequence Noise with integer CV
steps. Keep pitch fixed at first, then vary it to move the noise clock.
For a coupled tone, enable High-Pass Tone 1 from Tone 3 and sweep Tone 3's
frequency while monitoring Tone 1. Tone 3 now affects the result even when
you are not listening to its separate output.

Changing global clock switches can shift pitches abruptly. Return them to
the initialized positions before comparing tuning. Do not infer the actual
frequency from the current hover label; use a tuner, especially at the
high end of the 8-bit period range. This guide does not claim calibrated
tracking across every clock mode.

Global switch CV uses Rack's mono-to-poly repetition. The per-voice pitch,
noise and level inputs use indexed channels as described below. No extra
context-menu setting or custom wave bank is stored with this module.

\subsection{Cable channels, outputs and saved patches}
The widest input selects 1--16 polyphonic chip instances. Voice columns are
oscillators within each instance, not separate cable channels. Inputs read
the matching channel; supply matching pitch and level channel counts rather
than expecting a mono envelope to repeat across a polyphonic pitch cable.

Each output adds the preceding output when that earlier socket is unpatched.
The accumulated signal clips at approximately $\pm5$\,V. To hear the complete
mix, use the last output and leave earlier outputs empty. Patching an earlier
output removes that signal from the following mix. Lower source levels if
the sum clips.

Rack saves the panel parameters. Reset initializes the chip; oscillator
phase is not preserved by a patch save. Pitch/FM update each sample, while
level and other register controls update every 16 samples.
